\chapter{Implementation}

This chapter presents the Oracle object-relational implementation derived from the final logical design and UML schema.

\section{Types Definition}


\begin{lstlisting}
CREATE OR REPLACE TYPE MunicipalityListTY AS TABLE OF VARCHAR2(80);
/
\end{lstlisting}

\begin{lstlisting}
CREATE OR REPLACE TYPE DepotTY AS OBJECT (
    ID            VARCHAR2(20),
    Name          VARCHAR2(100),
    Address       VARCHAR2(100),
    City          VARCHAR2(50),
    Province      VARCHAR2(50),
    N_Emp         NUMBER,
    CoveredRegion VARCHAR2(80),
    Municipalities MunicipalityListTY
);
/
\end{lstlisting}

\begin{lstlisting}
CREATE OR REPLACE TYPE MemberTY AS OBJECT (
    SSN     VARCHAR2(16),
    Name    VARCHAR2(50),
    Surname VARCHAR2(50)
);
/

CREATE OR REPLACE TYPE MemberListTY AS TABLE OF MemberTY;
/

CREATE OR REPLACE TYPE SetupTeamTY AS OBJECT (
    Code    VARCHAR2(20),
    Name    VARCHAR2(100),
    Depot   REF DepotTY,
    Members MemberListTY
);
/
\end{lstlisting}

\begin{lstlisting}
CREATE OR REPLACE TYPE CustomerTY AS OBJECT (
    TaxCode       VARCHAR2(20),
    CustomerType  VARCHAR2(20),
    FirstName     VARCHAR2(50),
    LastName      VARCHAR2(50),
    VAT           VARCHAR2(20),
    CompanyName   VARCHAR2(100)
);
/
\end{lstlisting}

\begin{lstlisting}
CREATE OR REPLACE TYPE EventLocationTY AS OBJECT (
    Code     VARCHAR2(20),
    Address  VARCHAR2(100),
    NumberH  NUMBER,
    City     VARCHAR2(50),
    Prov     VARCHAR2(50),
    ZIP      VARCHAR2(10),
    Customer REF CustomerTY
);
/
\end{lstlisting}

\begin{lstlisting}
CREATE OR REPLACE TYPE BookingTY AS OBJECT (
    Code          VARCHAR2(20),
    Type          CHAR(1),
    Cost          NUMBER(10,2),
    BookDate      DATE,
    IntervalM     NUMBER,
    IsSpecial     CHAR(1),
    Team          REF SetupTeamTY,
    EventLocation REF EventLocationTY
);
/
\end{lstlisting}


\section{Tables Definition}


\begin{lstlisting}
CREATE TABLE DEPOT OF DepotTY (
    ID PRIMARY KEY,
    Name NOT NULL,
    Address NOT NULL,
    City NOT NULL,
    Province NOT NULL,
    N_Emp CHECK (N_Emp >= 0)
)
NESTED TABLE Municipalities STORE AS DepotMunicipalitiesNT;
\end{lstlisting}

\begin{lstlisting}
CREATE TABLE SETUPTEAM OF SetupTeamTY (
    Code PRIMARY KEY,
    Name NOT NULL,
    Depot NOT NULL REFERENCES DEPOT
)
NESTED TABLE Members STORE AS SetupTeamMembersNT;
\end{lstlisting}

\begin{lstlisting}
CREATE TABLE CUSTOMER OF CustomerTY (
    TaxCode PRIMARY KEY,
    CustomerType NOT NULL CHECK (CustomerType IN ('individual', 'company')),
    VAT UNIQUE,
    CHECK (
        (CustomerType = 'individual'
            AND FirstName IS NOT NULL
            AND LastName IS NOT NULL
            AND VAT IS NULL
            AND CompanyName IS NULL)
        OR
        (CustomerType = 'company'
            AND VAT IS NOT NULL
            AND CompanyName IS NOT NULL
            AND FirstName IS NULL
            AND LastName IS NULL)
    )
);
\end{lstlisting}

\begin{lstlisting}
CREATE TABLE EVENTLOCATION OF EventLocationTY (
    Code PRIMARY KEY,
    Address NOT NULL,
    NumberH CHECK (NumberH > 0),
    City NOT NULL,
    Prov NOT NULL,
    ZIP NOT NULL,
    Customer NOT NULL REFERENCES CUSTOMER ON DELETE CASCADE
);
\end{lstlisting}

\begin{lstlisting}
CREATE TABLE BOOKING OF BookingTY (
    Code PRIMARY KEY,
    Type NOT NULL CHECK (Type IN ('P','M','E','W')),
    Cost NOT NULL CHECK (Cost >= 0),
    BookDate NOT NULL,
    IntervalM NOT NULL CHECK (IntervalM >= 0),
    IsSpecial NOT NULL CHECK (IsSpecial IN ('Y','N')),
    Team NOT NULL REFERENCES SETUPTEAM,
    EventLocation NOT NULL REFERENCES EVENTLOCATION
);
\end{lstlisting}

\section{Trigger Implementation}


\subsection*{CheckCustomerType}
Enforces total and exclusive specialization for customers.

\begin{lstlisting}
CREATE OR REPLACE TRIGGER trg_check_customer_type
BEFORE INSERT OR UPDATE ON CUSTOMER
FOR EACH ROW
BEGIN
    IF :NEW.CustomerType = 'individual' THEN
        IF :NEW.FirstName IS NULL OR :NEW.LastName IS NULL THEN
            RAISE_APPLICATION_ERROR(-20001, 'Individual customer requires first and last name');
        END IF;
        IF :NEW.VAT IS NOT NULL OR :NEW.CompanyName IS NOT NULL THEN
            RAISE_APPLICATION_ERROR(-20002, 'Individual customer cannot have company data');
        END IF;
    ELSIF :NEW.CustomerType = 'company' THEN
        IF :NEW.VAT IS NULL OR :NEW.CompanyName IS NULL THEN
            RAISE_APPLICATION_ERROR(-20003, 'Company customer requires VAT and company name');
        END IF;
        IF :NEW.FirstName IS NOT NULL OR :NEW.LastName IS NOT NULL THEN
            RAISE_APPLICATION_ERROR(-20004, 'Company customer cannot have personal name data');
        END IF;
    ELSE
        RAISE_APPLICATION_ERROR(-20005, 'Invalid customer type');
    END IF;
END;
/
\end{lstlisting}

\subsection*{CheckDepotMunicipalities}
Guarantees that each depot covers at least one municipality.

\begin{lstlisting}
CREATE OR REPLACE TRIGGER trg_check_depot_municipalities
BEFORE INSERT OR UPDATE ON DEPOT
FOR EACH ROW
BEGIN
    IF :NEW.Municipalities IS NULL OR :NEW.Municipalities.COUNT = 0 THEN
        RAISE_APPLICATION_ERROR(-20006, 'Depot must include at least one municipality');
    END IF;
END;
/
\end{lstlisting}

\subsection*{CheckTeamMembers}
Enforces composition cardinality \texttt{SetupTeam 1 *-- 0..10 Member} and uniqueness of member SSN within each team.

\begin{lstlisting}
CREATE OR REPLACE TRIGGER trg_check_team_members
BEFORE INSERT OR UPDATE ON SETUPTEAM
FOR EACH ROW
DECLARE
    v_dup NUMBER;
BEGIN
    IF :NEW.Members IS NOT NULL AND :NEW.Members.COUNT > 10 THEN
        RAISE_APPLICATION_ERROR(-20007, 'A setup team cannot contain more than 10 members');
    END IF;

    IF :NEW.Members IS NOT NULL THEN
        SELECT COUNT(*)
        INTO v_dup
        FROM (
            SELECT m.SSN
            FROM TABLE(:NEW.Members) m
            GROUP BY m.SSN
            HAVING COUNT(*) > 1
        );

        IF v_dup > 0 THEN
            RAISE_APPLICATION_ERROR(-20008, 'Duplicate member SSN in the same setup team');
        END IF;
    END IF;
END;
/
\end{lstlisting}

\subsection*{CheckBookingOverlap}
Prevents overlapping booking windows for the same setup team.

\begin{lstlisting}
CREATE OR REPLACE TRIGGER trg_check_booking_overlap
BEFORE INSERT OR UPDATE ON BOOKING
FOR EACH ROW
DECLARE
    v_overlap NUMBER;
BEGIN
    SELECT COUNT(*)
    INTO v_overlap
    FROM BOOKING b
    WHERE b.Team = :NEW.Team
      AND b.Code <> NVL(:OLD.Code, '###')
      AND :NEW.BookDate <= b.BookDate + b.IntervalM
      AND b.BookDate <= :NEW.BookDate + :NEW.IntervalM;

    IF v_overlap > 0 THEN
        RAISE_APPLICATION_ERROR(-20009, 'Team already assigned in an overlapping time window');
    END IF;
END;
/
\end{lstlisting}

\subsection*{CheckBookingEventLocationCustomer}
Ensures that each booking references an event location that is effectively linked to an existing customer profile.

\begin{lstlisting}
CREATE OR REPLACE TRIGGER trg_booking_eventlocation_customer_match
BEFORE INSERT OR UPDATE ON BOOKING
FOR EACH ROW
DECLARE
    v_location EventLocationTY;
    v_customer CustomerTY;
BEGIN
    IF :NEW.EventLocation IS NULL THEN
        RAISE_APPLICATION_ERROR(-20010, 'Booking must reference an event location');
    END IF;

    SELECT DEREF(:NEW.EventLocation)
    INTO v_location
    FROM DUAL;

    IF v_location IS NULL OR v_location.Customer IS NULL THEN
        RAISE_APPLICATION_ERROR(-20011, 'Event location must be linked to a customer');
    END IF;

    SELECT DEREF(v_location.Customer)
    INTO v_customer
    FROM DUAL;

    IF v_customer IS NULL THEN
        RAISE_APPLICATION_ERROR(-20012, 'Referenced customer does not exist');
    END IF;
END;
/
\end{lstlisting}

\subsection*{CheckBookingDateNotPast}
Blocks inserts and updates when the booking date is in the past.

\begin{lstlisting}
CREATE OR REPLACE TRIGGER trg_booking_date_not_past
BEFORE INSERT OR UPDATE ON BOOKING
FOR EACH ROW
BEGIN
    IF :NEW.BookDate < TRUNC(SYSDATE) THEN
        RAISE_APPLICATION_ERROR(-20013, 'Booking date cannot be in the past');
    END IF;
END;
/
\end{lstlisting}

\subsection*{LockPastBookingReassignment}
Prevents team or destination location reassignment for historical bookings.

\begin{lstlisting}
CREATE OR REPLACE TRIGGER trg_lock_past_booking_reassignment
BEFORE UPDATE OF Team, EventLocation ON BOOKING
FOR EACH ROW
BEGIN
    IF :OLD.BookDate < TRUNC(SYSDATE)
       AND (:OLD.Team <> :NEW.Team OR :OLD.EventLocation <> :NEW.EventLocation) THEN
        RAISE_APPLICATION_ERROR(-20014, 'Team or location cannot be changed for past bookings');
    END IF;
END;
/
\end{lstlisting}

\section{Database Population}

Data generation is organized into modular procedures, one per logical area, plus a final orchestration procedure.

\subsection*{populateDepot}

\begin{lstlisting}
CREATE OR REPLACE PROCEDURE populateDepot(p_num_depots IN NUMBER) AS
    v_emp_count NUMBER;
BEGIN
    FOR d IN 1 .. p_num_depots LOOP
        v_emp_count := TRUNC(DBMS_RANDOM.VALUE(8, 60));

        INSERT INTO DEPOT VALUES (
            DepotTY(
                'D' || TO_CHAR(d, 'FM000'),
                'Depot ' || DBMS_RANDOM.STRING('U', 6),
                'Street ' || DBMS_RANDOM.STRING('U', 8),
                'City ' || DBMS_RANDOM.STRING('U', 6),
                'PR' || TO_CHAR(MOD(d, 20), 'FM00'),
                v_emp_count,
                'Region ' || DBMS_RANDOM.STRING('U', 5),
                MunicipalityListTY(
                    'Municipality ' || DBMS_RANDOM.STRING('U', 6),
                    'Municipality ' || DBMS_RANDOM.STRING('U', 6),
                    'Municipality ' || DBMS_RANDOM.STRING('U', 6)
                )
            )
        );
    END LOOP;
END;
/
\end{lstlisting}

\subsection*{populateSetupTeam}

\begin{lstlisting}
CREATE OR REPLACE PROCEDURE populateSetupTeam(p_num_teams_per_depot IN NUMBER) AS
    v_team_code    VARCHAR2(20);
    v_members      MemberListTY;
    v_member_count NUMBER;
BEGIN
    FOR d IN (SELECT ID FROM DEPOT ORDER BY ID) LOOP
        FOR t IN 1 .. p_num_teams_per_depot LOOP
            v_team_code := 'T' || SUBSTR(d.ID, 2) || '_' || TO_CHAR(t, 'FM00');
            v_members := MemberListTY();
            v_member_count := TRUNC(DBMS_RANDOM.VALUE(0, 11));

            FOR m IN 1 .. v_member_count LOOP
                v_members.EXTEND;
                v_members(v_members.COUNT) := MemberTY(
                    'SSN' || TO_CHAR(TO_NUMBER(SUBSTR(d.ID, 2)) * 10000 + t * 100 + m, 'FM0000000'),
                    'Name' || DBMS_RANDOM.STRING('U', 7),
                    'Surname' || DBMS_RANDOM.STRING('U', 7)
                );
            END LOOP;

            INSERT INTO SETUPTEAM VALUES (
                SetupTeamTY(
                    v_team_code,
                    'Team ' || DBMS_RANDOM.STRING('U', 6),
                    (SELECT REF(dp) FROM DEPOT dp WHERE dp.ID = d.ID),
                    v_members
                )
            );
        END LOOP;
    END LOOP;
END;
/
\end{lstlisting}

\subsection*{populateCustomer}

\begin{lstlisting}
CREATE OR REPLACE PROCEDURE populateCustomer(p_num_customers IN NUMBER) AS
BEGIN
    FOR c IN 1 .. p_num_customers LOOP
        IF DBMS_RANDOM.VALUE(0, 1) < 0.6 THEN
            INSERT INTO CUSTOMER VALUES (
                CustomerTY(
                    'C' || TO_CHAR(c, 'FM0000'),
                    'individual',
                    'Name' || DBMS_RANDOM.STRING('U', 8),
                    'Surname' || DBMS_RANDOM.STRING('U', 8),
                    NULL,
                    NULL
                )
            );
        ELSE
            INSERT INTO CUSTOMER VALUES (
                CustomerTY(
                    'C' || TO_CHAR(c, 'FM0000'),
                    'company',
                    NULL,
                    NULL,
                    'VAT' || TO_CHAR(c, 'FM000000000'),
                    'Company ' || DBMS_RANDOM.STRING('U', 8)
                )
            );
        END IF;
    END LOOP;
END;
/
\end{lstlisting}

\subsection*{populateEventLocation}

\begin{lstlisting}
CREATE OR REPLACE PROCEDURE populateEventLocation(p_num_locations_per_customer IN NUMBER) AS
    v_location_code VARCHAR2(20);
BEGIN
    FOR c IN (SELECT TaxCode FROM CUSTOMER ORDER BY TaxCode) LOOP
        FOR l IN 1 .. p_num_locations_per_customer LOOP
            v_location_code := 'L' || SUBSTR(c.TaxCode, 2) || '_' || TO_CHAR(l, 'FM00');

            INSERT INTO EVENTLOCATION VALUES (
                EventLocationTY(
                    v_location_code,
                    'Address ' || DBMS_RANDOM.STRING('U', 10),
                    TRUNC(DBMS_RANDOM.VALUE(1, 300)),
                    'City ' || DBMS_RANDOM.STRING('U', 6),
                    'PR' || TO_CHAR(MOD(TO_NUMBER(SUBSTR(c.TaxCode, 2)), 20), 'FM00'),
                    'ZIP' || TO_CHAR(TRUNC(DBMS_RANDOM.VALUE(10000, 99999))),
                    (SELECT REF(cu) FROM CUSTOMER cu WHERE cu.TaxCode = c.TaxCode)
                )
            );
        END LOOP;
    END LOOP;
END;
/
\end{lstlisting}

\subsection*{populateBooking}

\begin{lstlisting}
CREATE OR REPLACE PROCEDURE populateBooking(p_num_bookings IN NUMBER) AS
    TYPE refTeamTab IS TABLE OF REF SetupTeamTY INDEX BY PLS_INTEGER;
    TYPE refLocTab  IS TABLE OF REF EventLocationTY INDEX BY PLS_INTEGER;
    TYPE numTab     IS TABLE OF NUMBER INDEX BY PLS_INTEGER;
    teams      refTeamTab;
    locations  refLocTab;
    v_last_end numTab;
    v_t_idx    PLS_INTEGER;
    v_l_idx    PLS_INTEGER;
    v_type     CHAR(1);
    v_interval NUMBER;
    v_start    NUMBER;
    v_rand     NUMBER;
BEGIN
    SELECT REF(t) BULK COLLECT INTO teams FROM SETUPTEAM t;
    SELECT REF(l) BULK COLLECT INTO locations FROM EVENTLOCATION l;

    FOR i IN 1 .. teams.COUNT LOOP
        v_last_end(i) := 0;
    END LOOP;

    FOR b IN 1 .. p_num_bookings LOOP
        v_t_idx := TRUNC(DBMS_RANDOM.VALUE(1, teams.COUNT + 1));
        v_l_idx := TRUNC(DBMS_RANDOM.VALUE(1, locations.COUNT + 1));

        v_rand := DBMS_RANDOM.VALUE(0, 1);
        IF v_rand < 0.25 THEN
            v_type := 'P';
        ELSIF v_rand < 0.50 THEN
            v_type := 'M';
        ELSIF v_rand < 0.75 THEN
            v_type := 'E';
        ELSE
            v_type := 'W';
        END IF;

        v_interval := TRUNC(DBMS_RANDOM.VALUE(0, 4));
        v_start := v_last_end(v_t_idx) + 1 + TRUNC(DBMS_RANDOM.VALUE(0, 3));
        v_last_end(v_t_idx) := v_start + v_interval;

        INSERT INTO BOOKING VALUES (
            BookingTY(
                'B' || TO_CHAR(b, 'FM000000'),
                v_type,
                ROUND(DBMS_RANDOM.VALUE(80, 1200), 2),
                TRUNC(SYSDATE) + v_start,
                v_interval,
                CASE WHEN DBMS_RANDOM.VALUE(0, 1) < 0.3 THEN 'Y' ELSE 'N' END,
                teams(v_t_idx),
                locations(v_l_idx)
            )
        );
    END LOOP;
END;
/
\end{lstlisting}

\begin{lstlisting}
CREATE OR REPLACE PROCEDURE PopulateDatabase(
    p_num_depots                  IN NUMBER,
    p_num_teams_per_depot         IN NUMBER,
    p_num_customers               IN NUMBER,
    p_num_locations_per_customer  IN NUMBER,
    p_num_bookings                IN NUMBER
)
AS
BEGIN
    populateDepot(p_num_depots);
    populateSetupTeam(p_num_teams_per_depot);
    populateCustomer(p_num_customers);
    populateEventLocation(p_num_locations_per_customer);
    populateBooking(p_num_bookings);
END;
/
\end{lstlisting}

\begin{lstlisting}
BEGIN
    DBMS_OUTPUT.PUT_LINE('Starting StageUp population...');
    PopulateDatabase(
        p_num_depots                 => 10,
        p_num_teams_per_depot        => 10,
        p_num_customers              => 500,
        p_num_locations_per_customer => 2,
        p_num_bookings               => 50000
    );
    DBMS_OUTPUT.PUT_LINE('StageUp population completed.');
END;
/
\end{lstlisting}

\section{Operations}
\subsection*{Operation 1: Register a new customer}

\begin{lstlisting}
CREATE OR REPLACE PROCEDURE proc_register_customer(
    p_taxcode       IN VARCHAR2,
    p_customer_type IN VARCHAR2,
    p_first_name    IN VARCHAR2,
    p_last_name     IN VARCHAR2,
    p_vat           IN VARCHAR2,
    p_company_name  IN VARCHAR2
) AS
BEGIN
    INSERT INTO CUSTOMER VALUES (
        CustomerTY(
            p_taxcode,
            p_customer_type,
            p_first_name,
            p_last_name,
            p_vat,
            p_company_name
        )
    );
END;
/
\end{lstlisting}

\subsection*{Operation 2: Record a new booking}

\begin{lstlisting}
CREATE OR REPLACE PROCEDURE proc_add_booking(
    p_code          IN VARCHAR2,
    p_type          IN CHAR,
    p_cost          IN NUMBER,
    p_book_date     IN DATE,
    p_interval      IN NUMBER,
    p_is_special    IN CHAR,
    p_team_code     IN VARCHAR2,
    p_location_code IN VARCHAR2
) AS
BEGIN
    INSERT INTO BOOKING VALUES (
        BookingTY(
            p_code,
            p_type,
            p_cost,
            p_book_date,
            p_interval,
            p_is_special,
            (SELECT REF(st) FROM SETUPTEAM st WHERE st.Code = p_team_code),
            (SELECT REF(el) FROM EVENTLOCATION el WHERE el.Code = p_location_code)
        )
    );
END;
/
\end{lstlisting}

\subsection*{Operation 3: Register a new event location}

\begin{lstlisting}
CREATE OR REPLACE PROCEDURE proc_add_event_location(
    p_code            IN VARCHAR2,
    p_address         IN VARCHAR2,
    p_number          IN NUMBER,
    p_city            IN VARCHAR2,
    p_prov            IN VARCHAR2,
    p_zip             IN VARCHAR2,
    p_customer_taxcode IN VARCHAR2
) AS
BEGIN
    INSERT INTO EVENTLOCATION VALUES (
        EventLocationTY(
            p_code,
            p_address,
            p_number,
            p_city,
            p_prov,
            p_zip,
            (SELECT REF(cu) FROM CUSTOMER cu WHERE cu.TaxCode = p_customer_taxcode)
        )
    );
END;
/
\end{lstlisting}

\subsection*{Operation 4: View team for a specific event location}

\begin{lstlisting}
CREATE OR REPLACE FUNCTION func_get_team_by_location(
    p_location_code IN VARCHAR2
) RETURN VARCHAR2 AS
    v_team VARCHAR2(200);
BEGIN
    SELECT st.Code || ' - ' || st.Name
    INTO v_team
    FROM BOOKING b
    JOIN SETUPTEAM st ON DEREF(b.Team).Code = st.Code
    WHERE DEREF(b.EventLocation).Code = p_location_code
    ORDER BY b.BookDate DESC
    FETCH FIRST 1 ROW ONLY;

    RETURN v_team;
EXCEPTION
    WHEN NO_DATA_FOUND THEN
        RETURN 'No team found for this event location';
END;
/
\end{lstlisting}

\subsection*{Operation 5: Ranked event locations by number of bookings}

\begin{lstlisting}
CREATE OR REPLACE PROCEDURE proc_ranked_locations(
    p_result OUT SYS_REFCURSOR
) AS
BEGIN
    OPEN p_result FOR
        SELECT
            el.Code AS LocationCode,
            COUNT(b.Code) AS NumBookings
        FROM EVENTLOCATION el
        LEFT JOIN BOOKING b
            ON DEREF(b.EventLocation).Code = el.Code
        GROUP BY el.Code
        ORDER BY NumBookings DESC, el.Code;
END;
/
\end{lstlisting}






